\documentclass[12pt,a4paper]{article}
\usepackage[T1]{fontenc}
\usepackage[utf8]{inputenc}
\usepackage[portuguese]{babel}
\usepackage{geometry}
\usepackage{indentfirst}
\usepackage{microtype}
\geometry{top=3cm, bottom=3cm, left=3.5cm, right=3cm}
\setlength{\parindent}{2em}
\setlength{\parskip}{0.4ex}

% [ 302 ] marca a quebra de página do original
\newcommand{\pg}[1]{\thinspace\textup{[\,#1\,]}\thinspace}
% enunciado de tese numerada que corre com o texto (segue-se travessão)
\newcommand{\tese}[2]{\bigskip\noindent\textbf{#1.\quad #2}\;\textemdash\;}
% enunciado de tese numerada em linha própria (seguem-se as subdivisões a), b)...)
\newcommand{\teseB}[2]{\bigskip\noindent\textbf{#1.\quad #2}\par}
% subdivisão a), b), c)... com rubrica em itálico
\newcommand{\sub}[2]{\medskip\noindent#1)\quad\emph{#2},\;\textemdash\;}
% subdivisão a), b), c)... sem rubrica
\newcommand{\subp}[1]{\medskip\noindent#1)\quad}

% ---------------------------------------------------------------------------
% Notas: por omissão saem no pé da página (\footnote normal).
% Com  \def\notascorridas{}  passado na linha de comando, cada nota passa a
% correr no corpo do texto, em corpo menor, num bloco recuado que interrompe
% o texto: a leitura só continua depois de a nota acabar.
%   latexmk -pdf "Artigo III.tex"                                  -> notas no pé
%   latexmk -pdf -jobname="Artigo III - notas no corpo" \
%           -pretex='\def\notascorridas{}' -usepretex "Artigo III.tex"  -> notas no corpo
% ---------------------------------------------------------------------------
\newif\ifnotascorridas
\ifdefined\notascorridas \notascorridastrue \fi
\ifnotascorridas
  \renewcommand\footnote[1]{%
    \stepcounter{footnote}\textsuperscript{\thefootnote}% marca no ponto do texto
    \par
    {\footnotesize
     \leftskip=2em \rightskip=2em \parindent=1em
     \addvspace{\medskipamount}%
     \noindent\textit{Nota \thefootnote.}\enspace #1\par
     \addvspace{\medskipamount}}%
  }
\fi

\begin{document}

\begin{center}
  {\large\textbf{Artigo III}}\\[0.4em]
  {\large\textbf{Essência e Existência}}
\end{center}

\vspace{1.5em}

\pg{302}

\tese{29}{Essência e existência} O \emph{ente}, como dissemos, denota uma coisa
que \emph{existe} ou pode \emph{existir}. Por isso o ente compõe-se de
\emph{essência} e de \emph{existência}, como de \emph{potência} e de
\emph{ato}. \textemdash{} Tendo declarado a ideia de \emph{ente} e a ideia de
\emph{potência} e de \emph{ato}, devemos agora tratar da \emph{essência} e da
\emph{existência}.\footnote{O assunto de que nos ocupamos neste artigo é de uma
importância suprema em toda a filosofia. Pedimos, pois, toda a atenção dos
nossos pacientes leitores.}

\tese{30}{Essência} \emph{Essência} é a coisa pela qual um \emph{ente} é
\emph{o que é}; ou, mais desenvolvidamente, é a coisa que constitui um
\emph{ente} numa certa \emph{espécie}, ou num certo \emph{gênero}, \pg{303} e o
distingue dos \emph{entes} das outras \emph{espécies}, ou dos outros
\emph{gêneros}. Assim a \emph{essência} de \emph{Pedro} é a \emph{humanidade},
porque é por \emph{ela} que \emph{Pedro} é \emph{homem} e não é outra coisa
\textemdash{} ou porque é \emph{ela} que constitui \emph{Pedro} na \emph{espécie
humana} e o distingue dos \emph{animais irracionais}.\footnote{São muitas e
várias as \emph{explicações} que se dão da \emph{essência}. Assim se diz também
que a \emph{essência} \textemdash{} é a coisa cujo \emph{ato} é \emph{o ser}, ou
a \emph{existência}; \textemdash{} é o que primeiramente se concebe num
\emph{ente}; \textemdash{} é o que faz conhecer a coisa como é em si mesma;
\textemdash{} é o que distingue uma coisa de todas as outras; \textemdash{} é a
raiz e o sujeito de todas as propriedades e operações de um \emph{ente}, etc.
Todavia, embora sejam muitas e várias as \emph{explicações} que se dão da
\emph{essência}, contudo todas exprimem uma coisa concebida como o \emph{fundo}
que constitui um ente numa certa categoria ou espécie e de que brotam, como do
seu princípio radical, todas as propriedades e operações do ente.

  Apresentamos \emph{explicações} da \emph{essência}, e não \emph{definições}.
  Porquanto a \emph{essência}, embora seja menos abstrata que o \emph{ente},
  contudo, pela sua simplicidade, não é suscetível de uma \emph{definição
  rigorosa} sem que se cometa uma \emph{petição de princípio}. Com efeito,
  devendo a \emph{definição rigorosa} explicar \emph{o que uma coisa é}, isto é,
  devendo explicar a \emph{essência} de uma coisa, não é possível formar uma
  \emph{definição rigorosa} sem que primeiramente se possua uma noção geral da
  própria \emph{essência}. A \emph{essência}, pois, é uma das noções primitivas
  e imediatamente evidentes, que não precisa de demonstração nem pode ser
  demonstrada.

  A \emph{essência} chama-se também \emph{quididade}, \emph{forma},
  \emph{natureza}; mas esses vocábulos denotam sempre uma e a mesma coisa,
  embora sob diferentes aspectos. \textemdash{} Chama-se \emph{quididade},
  porque a essência, constituindo as coisas no próprio gênero ou na própria
  espécie, é denotada pela definição, com a qual respondemos à pergunta:
  \emph{o que é isto \textemdash{} quid est res}? \textemdash{} Chama-se
  \emph{forma}, porque, assim como a forma é \emph{ato} e \emph{perfeição} da
  \emph{matéria}, porque a tira da indiferença e a constitui num determinado
  grau de ser, assim a essência é \emph{ato} e \emph{perfeição} do \emph{ente},
  que pode considerar-se como um sujeito constituído pela própria essência num
  determinado grau de ser. \textemdash{} Chama-se \emph{natureza} enquanto é o
  primeiro princípio da operação própria, porque nenhuma coisa é destituída da
  sua operação. \textemdash{} Notamos que a \emph{essência} se chama
  \emph{essência} enquanto é \emph{nela} e \emph{por ela} que uma coisa tem o
  \emph{ser} (\emph{esse}). Uma coisa tem o \emph{ser} (\emph{esse})
  \emph{na essência}, porque a essência é a \emph{potência} em que o \emph{ser}
  é recebido. Uma coisa tem o \emph{ser} (\emph{esse}) \emph{pela essência},
  porque é por causa da \emph{essência} que uma coisa tem o \emph{ser} pelo qual
  é constituída numa certa espécie ou gênero, e porque nenhuma coisa poderia
  receber o ser se não tivesse a capacidade ou a potência para isso, e a
  essência é essa capacidade. \pg{304} \textemdash{} Diz S. Tomás
  (\emph{De Ente et Essentia}, c.~1): «\textit{Illud per quod res constituitur
  in proprio genere vel specie, est quod significamus per definitionem
  indicantem quid est res; inde est quod nomen essentiae a philosophis in nomen
  quidditatis mutatur\dots{} Dicitur etiam forma, secundum quod per formam
  significatur perfectio vel certitudo uniuscuiusque rei\dots{} Nomen autem
  naturae videtur significare essentiam rei, secundum quod habet ordinem vel
  ordinationem ad propriam operationem rei, cum nulla res propria destituatur
  operatione\dots{} Sed essentia dicitur secundum quod per eam et in ea res
  habet esse}».

  Como se vê, a ideia de essência sucede imediatamente à de ente; porque, tendo
  conhecido um objeto sob a ideia indeterminada de ente, a primeira pergunta
  que fazemos acerca do mesmo objeto é a seguinte: o que é isto? \textemdash{}
  A resposta exprime a essência.}

\tese{31}{Divisão da essência} A \emph{essência} é \textemdash{} a)~\emph{possível}
e \emph{atual}; \textemdash{} b)~\emph{substancial} e \emph{acidental};
\textemdash{} c)~\emph{física} e \emph{metafísica}.

\sub{a}{Essência possível e atual} \emph{Possível} é a que presentemente não
existe, mas pode existir na ordem real. \textemdash{} \emph{Atual} é a que
presentemente existe na realidade. \textemdash{} Neste artigo tratamos da
essência \emph{atual}.

\sub{b}{Essência substancial e acidental} A \emph{substancial} é a que se
encontra na \emph{substância}; tal é a essência do \emph{homem}. \textemdash{} A
\emph{acidental} é a que se encontra no \emph{acidente}; tal é a essência da
\emph{sabedoria}. \textemdash{} A essência existe na \emph{substância} por um
modo \emph{absoluto}, no \emph{acidente} por um modo \emph{relativo}.

\sub{c}{Essência física e metafísica} É \emph{física} quando os seus elementos
são \emph{concretos}, \emph{reais} e \emph{realmente distintos}; assim a
essência \emph{física} do \emph{homem} é constituída pela \emph{alma racional} e
pelo \emph{corpo orgânico}. \textemdash{} É \emph{metafísica} quando os seus
elementos são \emph{abstratos}, \emph{lógicos} e \emph{logicamente distintos},
denotando o \emph{gênero} e a \emph{diferença}; assim a essência
\emph{metafísica} do \emph{homem} é constituída pela \emph{animalidade} e pela
\emph{racionalidade}.\footnote{A essência física é \emph{artificial} ou
\emph{natural}. É \emph{artificial} se a sua constituição depende do artista;
tal é a essência de um relógio. É \emph{natural} quando a sua constituição é
obra da natureza; tal é a essência de uma planta. \textemdash{} A essência
metafísica é apenas o conceito abstrato da própria essência física, dividido em
dois outros conceitos, que são o \emph{gênero} e a \emph{diferença}. Como
dissemos na Lógica, o gênero não exprime uma parte física ou realmente distinta
das outras, mas exprime a essência completa, ainda que de um modo
indeterminado. Assim, na definição do homem, o gênero \emph{animal} não denota
exclusivamente o corpo, mas denota o corpo e a alma, enquanto à perfeição e às
forças da vida vegetativa e sensitiva, abstraindo da perfeição e das forças da
vida intelectual, contidas no conceito de \emph{racional}, que exerce o múnus de
\emph{diferença} (cf.\ \emph{De Ente et Essentia}, c.~VI). \textemdash{} Mas as
partes, quer sejam físicas, quer sejam metafísicas, sempre se referem à
essência real; pois a essência física e a metafísica são uma e a mesma essência
real, considerada sob um duplo aspecto \textemdash{} nas coisas e na
inteligência.}

\teseB{32}{A essência é uma coisa real e objetiva}

\subp{a} A \emph{essência} é o que constitui o \emph{ente} na própria
\emph{espécie}, ou no próprio \emph{gênero}, e o distingue dos \emph{entes} das
outras \emph{espécies}, ou dos outros \emph{gêneros}. Ora o que constitui o
ente na própria \pg{305} espécie, ou no próprio gênero, e o distingue dos entes
das outras espécies, ou dos outros gêneros, não é um mero \emph{produto da
nossa fantasia}, mas é uma coisa \emph{real e objetiva}, porque as coisas são o
que são independentemente dos nossos \emph{atos}. Logo a \emph{essência} é uma
coisa \emph{real e objetiva}.

\subp{b} Toda a \emph{essência} é \emph{física} ou \emph{metafísica}. Se é
\emph{física}, é certamente \emph{real e objetiva}; assim o \emph{corpo
orgânico} e a \emph{alma racional} são coisas \emph{reais e objetivas}, e por
isso \emph{real e objetiva} deve ser a essência que é constituída por esses dois
elementos. \textemdash{} Se é \emph{metafísica}, é também \emph{real e
objetiva}; porque a essência \emph{metafísica}, sendo a noção \emph{abstrata} da
essência \emph{física}, é uma coisa que \textemdash{} ou existe na realidade,
embora por um modo diferente do modo por que existe na inteligência,
\textemdash{} ou na realidade tem o seu fundamento, e por isso é sempre uma
coisa \emph{real e objetiva}. Logo a \emph{essência} é uma coisa \emph{real e
objetiva}.\footnote{Veja-se o que dissemos na Lógica com relação à
objetividade das ideias universais. \textemdash{} São os \emph{sensualistas} e
os \emph{positivistas} que reduzem a essência das coisas à mera ficção da nossa
fantasia, a nomes ou conceitos exclusivamente subjetivos, aos quais nenhuma
realidade corresponde.}

\tese{33}{A essência encontra-se não só na substância, mas também no acidente}
A \emph{essência} é a coisa pela qual o \emph{ente} é \emph{o que é},
\textemdash{} ou que constitui o \emph{ente} numa certa \emph{espécie}, ou num
certo \emph{gênero}, e o distingue dos \emph{entes} das outras \emph{espécies},
ou dos outros \emph{gêneros}. Ora não só a \emph{substância}, mas também o
\emph{acidente} é o que é, pertence a uma certa \emph{espécie}, ou a um certo
\emph{gênero}, e distingue-se dos \emph{entes} das outras \emph{espécies}, ou
dos outros \emph{gêneros}. Logo a \emph{essência} encontra-se não só na
\emph{substância}, mas também no \emph{acidente}.\footnote{A conclusão é certa.
A \emph{essência} responde à pergunta: \emph{o que é isto}? Ora esta pergunta
pode fazer-se não só com relação à \emph{substância}, mas também com relação ao
\emph{acidente}. Porquanto, não só a \emph{substância}, p.~ex., o homem, é o
que é e distingue-se das plantas, das pedras, etc., mas também o
\emph{acidente}, p.~ex., a alvura, é o que é e distingue-se das outras cores, da
sabedoria, do poder, etc. \pg{306} \textemdash{} Todavia, a \emph{essência}
compete primeiramente à \emph{substância} e secundariamente ao \emph{acidente}.
Na verdade, a essência deriva do \emph{ser} (\emph{esse}), que existe nela e por
ela. Ora o ser primeiramente compete à \emph{substância} e secundariamente ao
\emph{acidente}, visto que o ser do \emph{acidente} supõe o ser da
\emph{substância} e neste se funda, isto é, o \emph{ser acidental} supõe o
\emph{ser substancial} e neste se apoia. Por isso, a \emph{essência
substancial} é uma forma pela qual um ente é constituído no seu \emph{ser
primeiro}, que é o substancial, e existe independentemente de toda e qualquer
forma anterior; tal é a \emph{humanidade}. Ao passo que a \emph{essência
acidental} é uma forma pela qual um ente, já constituído no seu primeiro ser, é
diversamente modificado no seu \emph{ser secundário}, que é o acidental; tal é
a beleza, a sabedoria, a bondade, etc. \textemdash{} Esta é a doutrina de
S. Tomás. Diz o Santo que a essência se deve atribuir à \emph{substância} e ao
\emph{acidente} como se atribui o \emph{ente}. Ora o ente atribui-se primeira e
absolutamente à \emph{substância}, e secundária e relativamente ao
\emph{acidente}. Logo também a essência deve atribuir-se primeira e
absolutamente à \emph{substância}, e secundária e relativamente ao
\emph{acidente} (cf.\ \emph{De Ente et Essentia}, c.~III).}

\teseB{34}{A essência substancial não é uma simples coleção de qualidades}

\subp{a} A \emph{essência substancial} denota a realidade que constitui o
\emph{ente} no seu \emph{ser primeiro} e o faz subsistir em si mesmo; ao passo
que a \emph{qualidade}, como todo e qualquer \emph{acidente}, denota uma coisa
que não constitui o ente no seu ser primeiro, mas sobrevém ao ente já
constituído e subsistente em si mesmo. Ora o que se une a um ente já
constituído no ser primeiro não pode entrar, como elemento, nessa primeira
constituição. Logo a \emph{essência substancial} não é uma simples coleção de
\emph{qualidades}.

\subp{b} Se a \emph{essência substancial} fosse apenas uma coleção de
\emph{qualidades}, quando a substância se muda nas suas \emph{qualidades}
deveria mudar-se também na sua \emph{essência}. Ora isto não sucede; porque,
enquanto as \emph{qualidades} mudam, aparecem e desaparecem, a \emph{essência}
permanece sempre a mesma; assim o homem de bom torna-se mau, mas é sempre
homem. Logo a \emph{essência substancial} não é uma simples coleção de
\emph{qualidades}.\footnote{A verdade da proposição resulta da própria
inspeção dos termos. A \emph{substância} subsiste em si e por isso, para
existir, não precisa de sujeito a que adira; ao passo que a \emph{qualidade},
por ser \emph{acidente}, não subsiste em si e por isso, para existir, precisa
de um sujeito a que adira. Não pode, pois, afirmar-se, sem contradição, que a
\emph{essência substancial} é apenas uma coleção de \emph{qualidades}.
\textemdash{} Nem se diga que, se uma \emph{qualidade} não pode subsistir em si
mesma, pode subsistir a coleção das \emph{qualidades}; porque a
\emph{não-subsistência} é o caráter essencial das \emph{qualidades}, e por
isso, se a \emph{subsistência} não se encontra numa \emph{qualidade}, nem se
encontra na coleção de \emph{qualidades}. Se um cego não vê, nem podem ver mil
cegos.

  Foi \emph{Locke}, seguido pelos \emph{materialistas}, quem ensinou
  \textemdash{} que a \emph{substância}, ou \emph{essência substancial}, não é
  uma coisa distinta das \emph{qualidades sensíveis}, \textemdash{} e que, por
  isso, uma \emph{substância} concebida como um princípio ou sujeito de que
  resultam ou a que aderem as \emph{qualidades} é uma ficção gratuita da nossa
  inteligência. \textemdash{} O erro de \emph{Locke} é uma consequência do seu
  \emph{sistema}, que admite a \emph{sensação} como única fonte de todos os
  conhecimentos humanos. \textemdash{} É falsa a afirmação de \emph{Locke},
  dizendo que, quando concebemos a \emph{essência substancial}, concebemo-la
  como um \emph{cúmulo} ou uma coleção de \emph{qualidades}; porque a nossa
  inteligência, quando, advertida pelos sentidos, percebe as \emph{qualidades},
  \pg{307} não pára aí, mas, através das \emph{qualidades} e por meio delas,
  chega até ao fundo do \emph{objeto} e descobre a \emph{essência} de que
  resultam e a que aderem as \emph{qualidades}.}

\tese{35}{Nem todas as essências nos são desconhecidas} A experiência atesta
que possuímos uma ideia clara e distinta de muitos seres; porque não só não
confundimos um com o outro, mas também indicamos os elementos que os constituem
numa certa e determinada \emph{espécie}. Ora não poderíamos extremar um ser de
outro, nem indicar os seus elementos constitutivos, se a \emph{essência} deles
nos fosse desconhecida; pois é pela \emph{essência} que um ser é o que é e se
distingue dos outros seres. Logo nem todas as \emph{essências} nos são
desconhecidas.\footnote{Dizemos que nem todas as \emph{essências} nos são
desconhecidas, porque concedemos que muitas escapam ao nosso alcance.
Acrescentamos que mesmo a noção que temos de algumas delas, embora seja clara e
distinta, é contudo inadequada e imperfeita; visto que a \emph{essência}, sendo
o que de mais íntimo se encontra nos entes, escapa a uma intuição direta e
imediata, e só é percebida pelo raciocínio, que, baseando-se no princípio de
\emph{causalidade}, sobe das propriedades e dos fenômenos que nos são
manifestos para o princípio de que derivam e que é a \emph{essência}.
\textemdash{} Mas, ainda que imperfeito, o conhecimento da essência é
\emph{científico}; porque o \emph{método} é científico quando se levanta, pela
indução, dos efeitos para a causa, das ações para a natureza do agente, dos
fenômenos para a essência.

  \emph{Locke}, \emph{Bonnet}, \emph{Condillac} dizem que o homem, se pode
  conhecer as \emph{essências nominais}, não pode conhecer as \emph{essências
  reais}. \textemdash{} \emph{Essências reais} são as que física e
  intrinsecamente constituem os entes, e por isso existem independentemente do
  nosso pensamento; enquanto as \emph{essências nominais} são as que a
  inteligência compõe por meio das ideias de gênero e de diferença específica,
  e por isso são meros produtos da imaginação, destituídos de todo o fundamento
  real.

  Esta opinião é falsa, não só porque sustenta que não podemos conhecer, de
  modo algum, nenhuma \emph{essência real}, mas também porque afirma que as
  \emph{essências nominais}, por serem compostas de gênero e de diferença
  específica, são meros produtos da nossa imaginação, sem fundamento real. O
  \emph{gênero} e a \emph{diferença} não são uma ficção inútil da nossa
  inteligência, mas são noções universais e abstratas, representativas de uma
  coisa que é real, ou que tem fundamento na realidade, como suficientemente
  explicamos. \pg{308}}

\tese{36}{A essência das substâncias corpóreas é composta} As \emph{substâncias
corpóreas}, de que é constituído o mundo, convêm entre si enquanto são
\emph{substâncias}, mas diferem enquanto uma \emph{substância} pertence a uma
\emph{espécie} e outra a outra \emph{espécie}. Ora esta \emph{semelhança} e esta
\emph{diferença} demonstram claramente que a \emph{essência} delas é
\emph{composta} de dois elementos: um \emph{comum}, de que deriva a
\emph{semelhança} e que é o \emph{gênero}; outro \emph{próprio}, de que deriva a
\emph{diferença} e que é a \emph{diferença específica}. Logo a \emph{essência}
das \emph{substâncias corpóreas} é \emph{composta}.\footnote{As
\emph{substâncias criadas} são \emph{corpóreas} e \emph{incorpóreas}. A
\emph{essência} das \emph{substâncias corpóreas} \textemdash{} enquanto existe
na realidade, é \emph{composta} de dois \emph{princípios}, ou de duas
\emph{causas intrínsecas}: uma \emph{material} e outra \emph{formal}, que
constituem um \emph{composto real}; \textemdash{} e, enquanto é conhecida e
definida pela nossa inteligência, resulta de dois \emph{conceitos}, um dos quais
é o \emph{conceito do gênero} e outro é o \emph{conceito de diferença
específica}, que constituem um \emph{composto intencional}. O \emph{gênero}
funda-se no \emph{princípio material}; a \emph{diferença}, no \emph{princípio
formal}. \textemdash{} A \emph{essência} das \emph{substâncias incorpóreas}
\textemdash{} enquanto existe fora da nossa inteligência, é uma \emph{forma
inteiramente simples}; \textemdash{} mas, enquanto é conhecida e definida por
nós, está sujeita a uma \emph{composição}. Porquanto a nossa inteligência, pela
sua imperfeição natural, apreende as coisas simples e espirituais pelo mesmo
modo por que apreende as \emph{essências materiais} e \emph{compostas}, e por
isso define as \emph{formas subsistentes} pelo \emph{gênero} e pela
\emph{diferença}, como define as \emph{essências materiais}. O \emph{gênero}
funda-se na \emph{potência}, a \emph{diferença} funda-se no \emph{ato}; porque
as \emph{substâncias simples}, embora não sejam \emph{compostas} quanto à
\emph{essência}, contudo são realmente \emph{compostas} de \emph{potência} e de
\emph{ato}, isto é, de \emph{essência} e de \emph{existência}.}

\tese{37}{Dotes das essências criadas} As \emph{essências criadas} são
\emph{necessárias}, \emph{indivisíveis} e \emph{eternas}.

\sub{a}{São necessárias, ou imutáveis} não com relação à \emph{existência}
(porque, neste sentido, são \emph{contingentes} e \emph{mutáveis}),
\textemdash{} mas com relação aos \emph{elementos específicos} que as
constituem; porque uma \emph{essência} não pode ser nem entender-se sem os
elementos de que é \emph{composta}. Assim a \emph{essência} do \emph{homem} é
\emph{necessária}, ou \emph{imutável}, com relação à \emph{animalidade} e à
\emph{racionalidade}; pois não pode conceber-se um \emph{homem} que não seja
\emph{animal} e \emph{racional}. Logo as \emph{essências criadas} são
\emph{necessárias}.\footnote{Toda a \emph{essência} pode considerar-se
\textemdash{} ou nos seus \emph{elementos específicos} e \emph{abstratos},
\textemdash{} ou nos seus \emph{elementos reais} e \emph{concretos}, conforme é
\emph{metafísica} ou \emph{física}. \textemdash{} Se a \emph{essência} se
considera nos seus \emph{elementos específicos}, isto é, na sua \emph{realidade
quiditativa}, \emph{abstraída} das condições que a individualizam ou determinam
neste ou naquele \emph{indivíduo}, importa no seu conceito uma coisa
\emph{absolutamente necessária e imutável}; porque reproduz em si mesma o
\emph{exemplar} que existe na inteligência de Deus e que é \emph{necessário} e
\emph{imutável}. \textemdash{} Se, porém, a \emph{essência} se considera nos
\emph{elementos reais} de que é \emph{composta} nos \emph{indivíduos
singulares}, nesse sentido não é \emph{necessária} nem \emph{imutável}, porque
pode ser alterada e corrompida pela \emph{ação dos agentes naturais}. Assim a
\emph{essência real} ou \emph{física} do \emph{homem} acaba com a morte, quando
a alma se separa do corpo. \textemdash{} Mas, quando a nossa inteligência
considera a \emph{essência abstraída} dos \emph{princípios individuais},
percebe unicamente a \emph{realidade quiditativa}, os \emph{elementos
específicos}.

  Como se vê, a \emph{necessidade} que atribuímos à \emph{essência metafísica}
  das \emph{criaturas} \textemdash{} não é \emph{absoluta}, porque Deus podia
  formar outros \emph{modelos} ou \pg{309} \emph{exemplares} que fossem
  reproduzidos na \emph{criação}, \textemdash{} mas é \emph{hipotética},
  enquanto, suposta a \emph{criação}, Deus não muda o seu plano, mas tudo
  executa em conformidade com a sua infinita \emph{Sabedoria}.}

\sub{b}{São indivisíveis} enquanto não podem receber ou perder um
\emph{elemento} sem que fiquem \emph{destruídas}. Tal \emph{indivisibilidade}
das \emph{essências} é uma consequência da sua \emph{necessidade}. De fato, se
a uma \emph{essência} se acrescentar ou tirar um \emph{elemento}, já não
permanece a mesma, deixa de ser o que é; assim, se à \emph{substância viva e
sensitiva}, que é o \emph{bruto}, se acrescenta \emph{racional}, não temos o
\emph{bruto}, mas \emph{homem}; e se à mesma \emph{substância} se tira o
\emph{elemento sensitivo}, não temos o \emph{bruto}, mas a \emph{planta}. Logo
as \emph{essências criadas}, embora \emph{compostas}, são
\emph{indivisíveis}.\footnote{As \emph{essências criadas}, tanto \emph{físicas}
como \emph{metafísicas}, por serem \emph{compostas}, são \emph{divisíveis} nas
partes respectivas. As \emph{essências físicas} das \emph{substâncias
corpóreas} são \emph{divisíveis} na \emph{matéria} e na \emph{forma
substancial}; e as \emph{essências metafísicas} são \emph{divisíveis}, pela
\emph{análise da razão}, nas partes da definição, que são o \emph{gênero} e a
\emph{diferença específica}. Não é, pois, neste sentido que as \emph{essências
criadas} são \emph{indivisíveis}. Portanto, as \emph{essências} podem
dizer-se \emph{indivisíveis} enquanto não admitem o \emph{mais} nem o
\emph{menos}, enquanto os seus elementos não são suscetíveis de \emph{adição}
ou de \emph{diminuição} sem que uma \emph{essência} deixe de ser o que é. Por
isso as \emph{essências} foram comparadas aos \emph{números}; porque, assim
como um \emph{número} muda de \emph{espécie} pela \emph{adição} ou
\emph{subtração} de uma \emph{unidade}, assim também as \emph{essências} mudam
de \emph{espécie} pela \emph{adição} ou \emph{subtração} de um \emph{elemento}
(\emph{Sum. Th.}, 1-2, q.~52, a.~1).}

\sub{c}{São eternas} não enquanto existem \emph{realmente} desde a
\emph{eternidade} (pois que, neste sentido, só Deus é \emph{eterno}),
\textemdash{} mas enquanto, consideradas \emph{em si} e \emph{em abstrato},
prescindem de toda e qualquer \emph{diferença de tempo}, e também enquanto
existem, desde a \emph{eternidade}, nas \emph{ideias arquétipas} do
\emph{Criador}. Assim foi, é e será sempre verdade que um \emph{triângulo} tem
\emph{três ângulos}. Logo as \emph{essências criadas} são
\emph{eternas}.\footnote{As \emph{essências criadas} podem considerar-se
\emph{em si mesmas} e na \emph{Inteligência divina}. \textemdash{} Consideradas
\emph{em si mesmas}, as \emph{essências} prescindem de toda a \emph{limitação}
do \emph{tempo} e do \emph{lugar}. É uma \emph{eternidade negativa}.
\textemdash{} Consideradas na \emph{Inteligência divina}, as \emph{essências
criadas} são dotadas de todos os atributos próprios da \emph{Divindade};
porque, nesta \emph{hipótese}, as \emph{essências criadas} não se distinguem
realmente da \emph{Inteligência divina}, nem, por isso, da \emph{Essência
divina}, mas são a própria \emph{Essência infinita}, enquanto pode ser
diversamente \emph{imitada} e \emph{participada} pelas \emph{criaturas}. As
\emph{essências criadas}, pois, são outras tantas \emph{participações}, ainda
que \emph{inadequadas}, da única e simplicíssima \emph{Essência infinita}. Esta
\emph{participação} não se faz por \emph{emanação}, mas por \emph{semelhança},
enquanto o \emph{efeito} é semelhante à sua \emph{causa}.}

\tese{38}{Existência} \emph{Existência} é o \emph{ato}, ou a \emph{realização}
da \emph{essência}; ou, mais claramente, é a \emph{perfeição} que tira a
\emph{essência}, em que é recebida, do \emph{estado de possibilidade},
\pg{310} colocando-a no \emph{estado de atualidade}. \textemdash{} Por isso, o
\emph{ente real} ou \emph{atual} é um \emph{composto de essência e de
existência}.\footnote{A \emph{existência}, como a \emph{essência}, é uma
\emph{noção simplicíssima} e por si evidente, e, à semelhança das outras noções
simples e evidentes, não é capaz nem precisa de uma \emph{definição rigorosa}.
Por isso a nossa não é senão uma \emph{explicação}. \textemdash{} A
\emph{existência} é e diz-se \emph{ato da essência}. O \emph{ato}, que
\textemdash{} no sentido vulgar significa \emph{operação}, \textemdash{} no
sentido rigoroso significa uma coisa que \emph{determina} ou \emph{aperfeiçoa}.
Ora a \emph{essência} é por si \emph{indiferente} quanto à \emph{existência},
pode recebê-la, e por isso é um sujeito \emph{indeterminado} e
\emph{determinável}, \emph{imperfeito} e \emph{perfectível}; e é a
\emph{existência} que, recebida nela, a tira da \emph{indiferença} ou
\emph{indeterminação}, dando-lhe o \emph{complemento} que lhe falta na linha de
\emph{ente} ou da \emph{atualidade}, e assim faz com que a \emph{essência}
subsista em si mesma, na realidade, e de \emph{possível} se torne
\emph{existente}. Antes de receber a \emph{existência}, a \emph{essência} era
\emph{possível}; recebida a \emph{existência}, a \emph{essência} é uma coisa
\emph{real} ou \emph{atual}, e opõe-se à \emph{essência possível}.

  A \emph{existência}, ou o \emph{ato de ser}, é o \emph{ato}, ou a
  \emph{atualidade} de todos os outros \emph{atos}, e por isso é a
  \emph{perfeição de todas as perfeições}. E nada mais exato. Porquanto todo o
  \emph{ato}, e por isso toda a \emph{perfeição} (que é ato), se é ato
  considerado em si mesmo, é contudo \emph{potência} considerado em relação ao
  \emph{ato de ser}; porque, nos \emph{entes finitos}, nenhum \emph{ato}, e por
  isso nenhuma \emph{perfeição}, encerra na sua \emph{essência} a
  \emph{existência}, ou o \emph{ser}, isto é, o ato pode existir, mas não existe
  \emph{essencialmente}. E, como o ato denota \emph{atualidade} e
  \emph{perfeição}, podemos e devemos dizer que o \emph{ato de ser}, ou a
  \emph{existência}, é a \emph{atualidade} de todos os outros \emph{atos}, e por
  isso é a \emph{perfeição de todas as perfeições}. \textemdash{} Ao \emph{ser}
  nada pode acrescentar-se que seja mais \emph{formal}, mais \emph{atual}, e que
  o aperfeiçoe e determine, como o \emph{ato} determina a \emph{potência}. Por
  isso o \emph{ser} não é determinado por outra coisa, como a \emph{potência} é
  determinada pelo \emph{ato}; mas só pode ser e é determinado pela
  \emph{potência}. E assim um \emph{ser} distingue-se de outro \emph{ser}
  enquanto um é recebido numa \emph{essência} ou \emph{natureza} e outro
  noutra (o \emph{ato} é determinado ou limitado pela \emph{potência} em que é
  recebido). Eis as palavras do S. Doutor: «\textit{Hoc quod dico esse est
  actualitas omnium actuum, et propter hoc est perfectio omnium perfectionum.
  Nec intelligendum est quod ei quod dico esse aliquid addatur quod sit eo
  formalius, ipsum determinans, sicut actus potentiam. Unde non sic
  determinatur esse per aliud, sicut potentia per actum, sed magis sicut actus
  per potentiam\dots{} Et per hunc modum hoc esse ab illo esse distinguitur, in
  quantum est talis vel talis naturae}» (\emph{De potentia}, q.~7, a.~2 ad~9;
  cf.\ \emph{Sum. Th.}, p.~1, q.~4, a.~1 ad~3). \textemdash{} Sendo a
  \emph{existência} o que de mais \emph{atual} se encontra nas coisas, retamente
  os \emph{escolásticos} definiram a \emph{existência}: \emph{a última
  atualidade de uma coisa}; porque uma coisa recebe a última \emph{atualidade},
  isto é, a última \emph{perfeição}, o último \emph{complemento}, quando recebe
  a \emph{existência}. \textemdash{} Por isso, se a \emph{essência} é boa, é boa
  não \emph{por si}, mas pelo \emph{ser}. O \emph{ser} é o fundamento e a razão
  de todo o bem, ou de toda a \emph{perfeição}, que se encontra nas coisas. Diz
  S. Tomás: «\textit{Omnis nobilitas cuiuscumque rei est sibi secundum suum
  esse; nulla enim nobilitas ex sua sapientia (id est, ex ipsa essentia
  sapientiae), nisi per eam sapiens esset, et sic de aliis. Sic ergo secundum
  modum quo res habent esse, est suus modus in nobilitate}» (\emph{C. Gent.},
  1, 98; cf.\ \emph{De Ver.}, q.~21, a.~1 ad~4; a.~5 ad~5; \emph{Sum. Th.},
  p.~1, q.~62, a.~3 ad~2 et~3). \pg{311}

  Como se vê, a \emph{ideia de existência} sucede imediatamente à de
  \emph{essência}; porque, depois de termos concebido o que uma coisa é
  (\emph{quid est}), logo em seguida perguntamos se ela existe (\emph{an est}).
  \textemdash{} Todavia, do fato de a \emph{existência} ser uma coisa distinta
  da \emph{essência}, não se segue que a \emph{existência} e a \emph{essência}
  sejam dois \emph{entes}. A \emph{essência}, distinta da \emph{existência}, não
  é; nem é a \emph{existência}, distinta da \emph{essência}; mas é o \emph{ente
  composto de essência e de existência}. \textemdash{} Mas não devemos imaginar
  que o composto de \emph{essência} e de \emph{existência} seja uma terceira
  coisa, resultante desses dois elementos e diversa de cada um deles, como o
  \emph{homem}, composto de \emph{alma} e de \emph{corpo}, é uma terceira coisa
  distinta de cada um desses dois elementos. Quando dizemos que a \emph{essência
  é}, o verbo \emph{é} denota a \emph{atualidade} da \emph{essência}, e a
  \emph{essência} é o sujeito dessa \emph{atualidade}, e assim temos a
  \emph{essência existente}, que não é uma terceira coisa distinta da
  \emph{existência} e da \emph{essência}, mas é a mesma \emph{essência} enquanto
  \emph{atuada} pelo \emph{ser} ou pela \emph{existência} (S. Tomás,
  \emph{Quodl.} III, a.~3 ad~1).

  Notamos que o \emph{ser} e a \emph{existência}, embora se empreguem como
  termos equivalentes, contudo não significam uma e a mesma coisa. Porquanto o
  \emph{ser} significa tudo o que um ente é e possui, modelado e vazado no
  molde da \emph{essência}, e por isso significa toda a \emph{perfeição} de que
  o ente é dotado, incluída a \emph{existência}; ao passo que a \emph{existência}
  por si indica apenas a \emph{posição do ente fora das suas causas}, e não tudo
  o que o ente possui.}

\tese{39}{A existência encontra-se na substância e no acidente} A
\emph{existência} é a \emph{última atualidade} de todo o \emph{ente real}. Ora
todo o \emph{ente real} que existe no mundo é \emph{substância} ou
\emph{acidente}. Logo a \emph{existência} encontra-se na \emph{substância} e no
\emph{acidente}. \textemdash{} Por isso, assim como há uma \emph{essência
substancial} e uma \emph{essência acidental}, assim também há uma
\emph{existência substancial} e uma \emph{existência acidental}. \textemdash{}
Todavia, a \emph{existência} principalmente convém à \emph{substância}, e
secundariamente ao \emph{acidente}.\footnote{A conclusão é evidente. No mesmo
indivíduo humano, uma coisa é a \emph{existência} ou o \emph{ser de homem}, e
outra coisa é a \emph{existência} ou o \emph{ser de sábio}, de bom, etc.
\textemdash{} A \emph{existência}, como a \emph{essência}, primeiramente compete
à \emph{substância} e secundariamente ao \emph{acidente}; porque o
\emph{acidente} tem o seu \emph{ser} ou a sua \emph{existência} no \emph{ser}
ou na \emph{existência} da \emph{substância}. \textemdash{} Em cada
\emph{substância} não pode haver senão um único \emph{ser substancial}, porque
não pode haver nela senão uma única \emph{forma substancial}, da qual deriva o
\emph{ser substancial}; mas pode haver vários \emph{seres acidentais}, porque
pode haver nela várias \emph{formas acidentais}, de cada uma das quais deriva um
\emph{ser acidental}. Assim \emph{Pedro} \textemdash{} que, em virtude da sua
única \emph{forma substancial}, isto é, da única \emph{alma racional}, possui o
\emph{ser substancial} pelo qual subsiste enquanto homem, \textemdash{} em
virtude de várias \emph{formas acidentais} possui vários \emph{seres
acidentais}, pelos quais é \emph{branco}, é \emph{virtuoso}, é \emph{sábio},
etc. (cf.\ Cajet., \emph{Com. in Op. De Ente et Essentia}, c.~7, q.~16).
\textemdash{} Advertimos que, nos entes compostos de \emph{matéria} e
\emph{forma}, o \emph{ser}, ou a \emph{existência}, se diz derivar da
\emph{forma} para o mesmo composto enquanto a \emph{forma}, aperfeiçoando e
completando a \emph{essência material} na linha da \emph{essência}, \pg{312} dá
à mesma \emph{essência} a \emph{última disposição}, necessária e suficiente para
a recepção da \emph{existência}, isto é, torna a mesma \emph{essência}
proximamente capaz de receber a \emph{existência}. Isto vale,
proporcionalmente, tanto para a \emph{forma substancial} quanto para a
\emph{forma acidental}. Da \emph{forma substancial}, que atua a \emph{matéria
prima}, deriva para o sujeito o \emph{ser substancial}; da \emph{forma
acidental}, que atua o sujeito já constituído no seu \emph{ser substancial}, que
é o \emph{primeiro ser}, deriva para ele o \emph{ser acidental}, que é o
\emph{segundo ser}.}

\tese{40}{Nos entes finitos, a existência refere-se à essência como o ato à
potência} Quando há duas coisas, uma das quais \emph{aperfeiçoa} e
\emph{completa} a outra, a primeira refere-se à segunda como o \emph{ato} se
refere à \emph{potência}; porque \emph{ato} é e diz-se a coisa que
\emph{aperfeiçoa} e \emph{completa}, e \emph{potência} é e diz-se a coisa que
recebe a \emph{perfeição} e o \emph{complemento}. Ora, nos \emph{entes finitos},
a \emph{existência} \emph{aperfeiçoa} e \emph{completa} a \emph{essência},
porque a faz subsistir na \emph{atualidade}. Logo, nos \emph{entes finitos}, a
\emph{existência} refere-se à \emph{essência} como o \emph{ato} à
\emph{potência}.\footnote{É o argumento de S. Tomás, que estabelece o seguinte
princípio: «\textit{In quocumque inveniuntur aliqua duo quorum unum est
complementum alterius, proportio unius ad alterum est sicut proportio potentiae
ad actum; nihil enim completur nisi per proprium actum}» (\emph{C. Gent.},
l.~II, c.~53). \textemdash{} E noutro lugar: «\textit{Quidditas est sicut
potentia, et suum esse acquisitum est sicut actus}» (\emph{In II Sent.},
dist.~3, q.~1, a.~1). \textemdash{} Sendo a \emph{existência} para a
\emph{essência} o que o \emph{ato} é para a \emph{potência}, tudo o que
dissemos com relação ao \emph{ato} e à \emph{potência} deve aplicar-se à
\emph{existência} e à \emph{essência}. \textemdash{} Donde se segue que a
\emph{existência} é \emph{simplesmente mais perfeita} do que a \emph{essência}.
Com efeito, a \emph{existência}, sendo a \emph{última atualidade} da
\emph{essência}, enquanto a faz passar do estado de \emph{ente em potência}
para o estado de \emph{ente em ato}, é para ela um princípio de
\emph{complemento} e de \emph{perfeição}. Ora o \emph{ente} que aperfeiçoa é
\emph{simplesmente} mais perfeito do que o \emph{ente} que é aperfeiçoado.

  Dissemos: \emph{simplesmente}, porque, sob um certo aspecto, a \emph{essência}
  é mais perfeita que a \emph{existência}. De fato, a \emph{essência}, por
  \emph{limitar} e \emph{determinar} a \emph{existência}, refere-se à própria
  \emph{existência} como o que limita e determina se refere ao que é limitado e
  determinado. Ora o \emph{determinante} é mais perfeito que o
  \emph{determinável}. \textemdash{} Além disso, o \emph{ser}, ou a
  \emph{existência}, embora não possa dizer-se \emph{simplesmente}
  aperfeiçoado pela \emph{essência} (a \emph{essência} é aperfeiçoada pela
  \emph{existência}, da qual recebe a \emph{última atualidade}), contudo,
  quando é recebido na \emph{essência}, é por esta de algum modo aperfeiçoado,
  enquanto se diz que a \emph{existência}, ou que o \emph{ser}, é mais ou menos
  nobre e perfeito conforme é mais ou menos nobre e perfeita a \emph{essência}
  em que é recebido; assim o \emph{ser da planta} é mais perfeito que o \emph{da
  pedra}, o \emph{ser do homem} é mais perfeito que o \emph{do animal}.
  \textemdash{} A razão disso é porque o \emph{primeiro Agente} produz o
  \emph{ser}, não como subsistente em si mesmo, mas em ordem a uma certa
  \emph{essência}, proporcionado a uma \emph{essência} (pois deve haver
  proporção entre a coisa recebida e o recipiente); e por isso o \emph{ser},
  embora receba de Deus, e não da \emph{essência}, a própria \emph{perfeição},
  contudo recebe-a de Deus \pg{313} em ordem à \emph{essência}; e Deus produz
  um \emph{ser} mais ou menos perfeito conforme a qualidade da \emph{essência}
  na qual o \emph{ser} é recebido como numa \emph{forma} ou \emph{molde}.
  Assim, a \emph{nobreza} e a \emph{perfeição} da \emph{essência} é o
  \emph{critério}, o \emph{sinal} da \emph{nobreza} e da \emph{perfeição} do
  \emph{ser}.}

\tese{41}{A existência reduz-se à categoria da essência} Tudo o que se refere a
uma \emph{essência}, como \emph{ato} à \emph{potência}, reduz-se à
\emph{categoria} em que se encontra a própria \emph{essência}; porque o
\emph{ente} coloca-se na \emph{categoria} em virtude da sua \emph{essência}, e a
\emph{potência} e o \emph{ato} pertencem ao mesmo \emph{gênero}, ou à mesma
\emph{categoria}. Ora a \emph{existência} refere-se à respectiva \emph{essência}
como \emph{ato} à \emph{potência}; porque a \emph{existência}, como dissemos, é
a \emph{última atualidade}, que completa a \emph{essência} e a coloca na
\emph{ordem real}. Logo a \emph{existência} reduz-se à \emph{categoria} da
\emph{essência}. \textemdash{} Por isso, a \emph{existência} da \emph{substância}
reduz-se à \emph{categoria} da mesma \emph{substância}, e é uma \emph{coisa
substancial}; e a \emph{existência} do \emph{acidente} reduz-se à
\emph{categoria} em que se encontra o mesmo \emph{acidente}, e é uma \emph{coisa
acidental}.\footnote{Referindo-se a \emph{existência} à \emph{essência} como o
\emph{ato} à \emph{potência}, poderia perguntar-se: a \emph{existência} será
apenas um \emph{acidente} da \emph{substância}, ou será uma \emph{parte
substancial}? \textemdash{} A esta pergunta damos a seguinte resposta, indicada
no texto. A \emph{existência} não é nem pode dizer-se \emph{acidente} da
\emph{substância}, se o \emph{acidente} se toma no sentido rigoroso, enquanto
significa uma coisa que, para existir, precisa de um sujeito a que adira. Na
verdade, a \emph{existência} dá a \emph{última atualidade} à \emph{substância},
colocando-a na \emph{ordem real}. Ora é absolutamente impossível que a
\emph{substância} receba do \emph{acidente} a \emph{última atualidade}.
\textemdash{} Mas, se o \emph{acidente} não se toma em sentido rigoroso, se por
\emph{acidente} se entende tudo o que não constitui a \emph{essência} de uma
coisa, neste sentido pode dizer-se que a \emph{existência} é \emph{acidente} da
\emph{substância}. Porquanto a \emph{existência} não entra na \emph{essência} de
uma coisa, não é um \emph{elemento constitutivo} da \emph{essência} (finita),
mas é a \emph{perfeição} que dá a \emph{última atualidade} à \emph{essência} já
constituída e completa na própria linha ou ordem da \emph{essência}, fazendo
com que a mesma \emph{essência} exista na \emph{ordem real}. Portanto devemos
dizer que a \emph{existência} é \emph{parte substancial} da \emph{substância},
ou é \emph{ato substancial} \textemdash{} não enquanto é um \emph{elemento
constitutivo} da própria \emph{substância}, ou da \emph{essência substancial},
\textemdash{} mas enquanto dá à \emph{substância} a \emph{última atualidade}, o
\emph{último complemento} na linha da realidade. Por isso, a \emph{existência}
não se reduz à \emph{categoria de substância} como a \emph{forma substancial}
se reduz à \emph{categoria do composto}, \textemdash{} mas como o \emph{ato} se
reduz à \emph{categoria da potência} (cf.\ \emph{De Pot.}, q.~5, a.~4 ad~3).

  Agora uma advertência muito importante. Dizemos que a \emph{existência}, sendo
  o \emph{complemento}, a \emph{última atualidade} da \emph{substância} e do
  \emph{acidente}, se reduz ao \emph{gênero de substância} ou de
  \emph{acidente}. Todavia, propriamente falando, a \emph{existência} não é nem
  \emph{substância} nem \emph{acidente categórico}; mas é uma coisa que
  \emph{transcende} toda a \emph{substância} e todo o \emph{acidente}, porque é
  \emph{ato do ente}, que não é \emph{gênero} e transcende todas as
  \emph{categorias}.} \pg{314}

\teseB{42}{O ente finito é composto de essência real e de existência real}

\subp{a} A \emph{essência} do \emph{ente finito}, considerada \emph{em si
mesma}, é \emph{necessária} e \emph{eterna}; e, todavia, considerada no mesmo
\emph{ente}, é \emph{contingente} e \emph{temporânea}. Ora a \emph{essência} do
\emph{ente finito} não é \emph{contingente} e \emph{temporânea} senão pela
\emph{existência}. Logo o \emph{ente finito} é \emph{composto de essência real e
de existência real}.

\subp{b} O \emph{ser} (ou a \emph{existência}), considerado \emph{em si mesmo},
é \emph{uno}, \emph{indiviso}, \emph{ilimitado}, \emph{não-causado}; e todavia,
considerado nos \emph{entes finitos}, é \emph{múltiplo}, \emph{diviso},
\emph{limitado}, \emph{causado}. Ora essa \emph{modificação} que o \emph{ser}
recebe nos \emph{entes finitos} só pode derivar da \emph{essência}, a qual,
sendo por si \emph{determinada}, é \emph{princípio de determinação}, e por isso
de \emph{divisão}, de \emph{multiplicidade}, de \emph{limitação} e de
\emph{dependência}. Logo o \emph{ente finito} é \emph{composto de existência
real e de essência real}.\footnote{O \emph{ser} é, por si, \emph{infinito} e
\emph{ilimitado}; porque a \emph{limitação} é o \emph{não-ser}. Portanto, a
\emph{limitação} de que o \emph{ser} se acha cercado nos \emph{entes finitos}
deriva da \emph{essência}, a qual, por ser \emph{determinada} e \emph{limitada}
em si mesma, \emph{determina} e \emph{limita} o \emph{ser} que ela recebe.
Assim, quando digo: \emph{o ser do homem}, não falo do \emph{ser em si}, do
\emph{ser puro}, que é só \emph{ser}, mas falo do \emph{ser próprio deste
ente}, e por isso de um \emph{ser limitado} e \emph{determinado}. \textemdash{}
Como também, o \emph{ser} é, por si, \emph{uno} e \emph{indiviso}. Se, pois, se
acha \emph{diviso} e \emph{multiplicado} nas \emph{criaturas}, essa
\emph{divisão} e \emph{multiplicação} deriva de outro \emph{princípio}, que é a
\emph{essência} (\emph{C. Gent.}, II, 52). Porquanto a \emph{multiplicação}
supõe a \emph{divisão}, e a \emph{divisão} importa uma \emph{oposição}. Ora tal
\emph{oposição} não pode dar-se da parte do próprio \emph{ser}, mas só da parte
das \emph{essências}, que, sendo por si \emph{determinadas}, são
\emph{distintas}, \emph{opostas}, \emph{múltiplas}, conforme \emph{imitam},
mais ou menos, a \emph{infinita e perfeitíssima Essência divina}. Diz S. Tomás:
«\textit{Determinatio est pluralitatis principium primum}» (\emph{Quodl.} VI,
a.~6). \textemdash{} Finalmente, o \emph{ser} das \emph{criaturas} deriva de
uma \emph{causa extrínseca}. Ora o \emph{ser} não é produzido \emph{por si},
enquanto é \emph{ser}. Diz S. Tomás: «\textit{Esse ab alio causatum non
competit enti inquantum est ens; alias omne ens esset ab alio causatum, et sic
oporteret procedere in infinitum in causis, quod est impossibile}» (l.~c.).
Logo, se o \emph{ser} das \emph{criaturas} não é produzido \emph{por si}, é
produzido por causa da \emph{essência} em que é recebido. \textemdash{} Logo o
\emph{ente finito} é \emph{composto} de um \emph{duplo princípio real}, ou
\emph{atual}, que é a \emph{essência} e a \emph{existência} (ou o \emph{ser}).

  Quando se diz que o \emph{ente finito} é \emph{composto} de \emph{essência
  real} e de \emph{existência} ou de \emph{ser real}, não deve entender-se que a
  \emph{realidade da essência} seja uma coisa distinta da \emph{realidade do
  ser}, de modo que o \emph{ente} resulte de \emph{duas realidades}. No
  \emph{ente finito} composto de \emph{essência} e de \emph{ser} há uma \emph{só
  realidade}. Esta \emph{única realidade} é própria do \emph{ser}, pertence ao
  \emph{ser}, e pelo \emph{ser} é comunicada à \emph{essência}; por forma que o
  \emph{ser} é \emph{real por si}, e a \emph{essência} se torna e é \emph{real}
  pela \emph{realidade do ser}, assim como a \emph{potência} se torna e é
  \emph{atual} pelo \emph{ato}, enquanto recebe em si o \emph{ato}.} \pg{315}

\tese{43}{No ente finito, a existência distingue-se logicamente da essência}
Uma coisa distingue-se \emph{logicamente} de outra quando o \emph{conceito} da
primeira é diverso do \emph{conceito} da segunda. Ora o \emph{conceito} da
\emph{existência} é diverso do \emph{conceito} da \emph{essência}; porque
\textemdash{} a \emph{essência} significa a coisa que constitui o \emph{ente}
numa determinada \emph{espécie}, ao passo que a \emph{existência} significa a
coisa que faz existir a mesma \emph{essência} na realidade; \textemdash{} a
\emph{essência} é capaz de \emph{existir}, mas é por si \emph{indeterminada}
quanto à sua presença na \emph{ordem real}, e por isso tem \emph{razão de
potência}, ao passo que a \emph{existência} tira da \emph{indeterminação} a
\emph{essência}, preenche aquela \emph{capacidade}, e por isso tem \emph{razão
de ato}. Logo, no \emph{ente finito}, a \emph{existência} distingue-se
\emph{logicamente} da \emph{essência}.\footnote{Que, em todo o \emph{ente
finito}, a \emph{existência} seja \emph{logicamente distinta} da \emph{essência},
é verdade admitida por todos os que sabem o que é \emph{distinção lógica} e o
que significam os termos \textemdash{} \emph{existência} e \emph{essência}.
\textemdash{} Como se vê, falamos da \emph{essência real}, ou \emph{atual}, de
que todas as \emph{criaturas} existentes são intrinsecamente \emph{compostas},
e não da \emph{essência possível}, que pode existir mas que atualmente não
existe na realidade.}

\teseB{44}{No ente finito, a existência distingue-se realmente da essência}

\subp{a} Como o nosso \emph{conhecimento} é uma \emph{reprodução ideal} das
coisas, se, no \emph{ente finito}, a \emph{existência} e a \emph{essência} não
se distinguissem \emph{realmente}, mas fossem uma e a mesma coisa, a
\emph{espécie} ou \emph{imagem inteligível} que nos representa a \emph{essência}
da coisa entendida deveria representar-nos também a sua \emph{existência},
porque uma coisa não pode separar-se de si mesma. Ora a \emph{experiência}
atesta que uma \emph{espécie} ou \emph{imagem inteligível} pode representar e
efetivamente representa a \emph{essência} das coisas, prescindindo contudo da
sua \emph{existência}. Logo, no \emph{ente finito}, a \emph{existência}
distingue-se \emph{realmente} da \emph{essência}.\footnote{É evidente. Se a
\emph{existência} fosse um \emph{elemento constitutivo} da \emph{essência}, toda
a \emph{definição} da \emph{essência} de um \emph{ente} deveria exprimir também
a sua \emph{existência}; visto que a \emph{definição} não pode deixar de
exprimir \emph{todos} os elementos que constituem a \emph{essência} do
\emph{ente} definido. Ora a \emph{definição} da \emph{essência} de um \emph{ente}
prescinde completamente da \emph{existência} ou da \emph{não-existência} do
mesmo \emph{ente}; porque podemos entender \emph{o que é o homem} e todavia
\emph{ignorar se existe na realidade}. Logo a \emph{existência}, no \emph{ente
finito}, não é um \emph{elemento constitutivo} da \emph{essência}, mas uma
distingue-se \emph{realmente} da outra. \textemdash{} Por outras palavras, se o
\emph{ser} da \emph{criatura} está \emph{fora da essência} quanto à
\emph{inteligência}, deve estar também \emph{fora da essência} quanto à
\emph{realidade}. \textemdash{} Diz S. Tomás: «\textit{Quidquid non est de
intellectu essentiae vel quidditatis, hoc est adveniens extra et faciens
compositionem cum essentia; quia \pg{316} nulla essentia sine his quae sunt
partes essentiae intelligi potest. Omnis autem essentia vel quidditas intelligi
potest sine hoc quod aliquid intelligatur de esse suo facto; possum enim
intelligere quid est homo vel phoenix, et tamen ignorare an esse habeant in
natura. Ergo patet quod esse est aliud ab essentia vel quidditate}» (\emph{De
Ente et Essentia}, c.~V). \textemdash{} Como se vê, S. Tomás fala claramente da
\emph{essência atual}, e não da \emph{possível}. Porquanto ensina ele que \emph{a
existência forma um composto com a essência}; ora a \emph{existência}, sendo
\emph{coisa real}, não pode formar o \emph{composto} com uma coisa que não
existe na realidade, como é a \emph{essência possível}.}

\subp{b} A \emph{existência} distinguir-se-á \emph{realmente} da \emph{essência}
se aquela não for um \emph{elemento constitutivo} desta, e por isso se a
\emph{essência} existir em virtude de uma outra \emph{realidade} ou
\emph{perfeição}; porque \emph{duas coisas são realmente distintas entre si
quando uma não é outra}. Ora, no \emph{ente finito}, a \emph{existência} não é
um \emph{elemento constitutivo} da \emph{essência}; porque, se o fosse, a
\emph{criatura} existiria em virtude da própria \emph{essência}, isto é, seria
\emph{essencialmente existente}, e por isso não seria \emph{contingente}, mas
\emph{necessária}, e assim deixaria de ser \emph{criatura}. Logo, no \emph{ente
finito}, a \emph{existência} distingue-se \emph{realmente} da
\emph{essência}.\footnote{Este argumento pode propor-se por este outro modo: se,
no \emph{ente finito}, a \emph{existência} não se distinguisse \emph{realmente}
da \emph{essência}, a \emph{existência} das \emph{criaturas} seria
\emph{necessária}, porque os elementos das \emph{essências} são
\emph{necessários}; e por isso, assim como o \emph{homem} é
\emph{necessariamente racional} (porque \emph{racional} é um \emph{elemento
constitutivo} da \emph{essência humana}), a \emph{criatura} seria
\emph{necessariamente existente}. Ora a \emph{existência} das \emph{criaturas}
não é \emph{necessária}, mas \emph{contingente}, porque podiam deixar de
existir. Logo a \emph{existência} não é um \emph{elemento constitutivo} da
\emph{essência} das \emph{criaturas}, e por isso, no \emph{ente finito}, a
\emph{existência} distingue-se \emph{realmente} da \emph{essência}.

  Pode objetar-se: \emph{uma coisa atualmente existente é-o necessariamente}.
  Logo não repugna que a \emph{essência atual} das \emph{criaturas} se diga e
  seja \emph{essencial} ou \emph{necessariamente existente}. Logo, no \emph{ente
  finito}, a \emph{existência} não se distingue \emph{realmente} da
  \emph{essência}.

  \emph{Resposta.} A \emph{necessidade na existência} pode referir-se ao
  \emph{fato} ou à \emph{exigência}. \textemdash{} Referida ao \emph{fato}, a
  \emph{necessidade na existência} é uma coisa \emph{hipotética}, e pode
  perfeitamente conciliar-se com a \emph{contingência} do \emph{ente}: de modo
  que um \emph{ente contingente}, enquanto existe, não pode deixar de existir,
  existe \emph{necessariamente}, e todavia não deixa de ser \emph{contingente},
  porque a \emph{existência} não lhe era \emph{necessária}, isto é, porque não
  existe por assim o exigir a sua \emph{essência}. \textemdash{} Referida à
  \emph{exigência}, a \emph{necessidade na existência} é uma coisa
  \emph{absoluta}, constitui a própria \emph{essência dos entes}, e por isso não
  pode conciliar-se com a \emph{contingência} dos próprios \emph{entes}: assim
  Deus, como veremos, existe por assim o exigir a sua \emph{essência}, e por
  isso é o \emph{ser absolutamente necessário}. \textemdash{} Ora a
  \emph{necessidade na existência} que se atribui aos \emph{entes finitos} ou
  \emph{criados} não é uma \emph{necessidade da exigência}, mas é uma
  \emph{necessidade de fato}; visto que os \emph{entes finitos} existem, não
  porque o exige a sua \emph{essência}, mas só porque Deus livremente os criou.
  Se os \emph{entes finitos} existem, não porque \pg{317} assim o exige a sua
  \emph{essência}, é claro que a sua \emph{necessidade na existência} não é um
  \emph{elemento constitutivo} da mesma \emph{essência}. Se, nos \emph{entes
  finitos}, a \emph{necessidade da existência} não é um \emph{elemento
  constitutivo} da \emph{essência}, a \emph{existência} distingue-se
  \emph{realmente} da \emph{essência}. \textemdash{} Nem pode dizer-se que a
  \emph{necessidade na existência} se torna um \emph{elemento constitutivo} da
  \emph{essência} dos \emph{entes finitos} pelo fato de estes existirem
  atualmente; porque, se assim fosse, as \emph{essências} não seriam
  \emph{imutáveis}, as \emph{criaturas} existiriam em virtude da sua
  \emph{essência} e deixariam de ser \emph{contingentes}: o que é
  \emph{absurdo}. Logo a \emph{essência} dos \emph{seres finitos} recebe ou tem
  a \emph{existência}, mas não é a \emph{existência}.}

\subp{c} Todo o \emph{ato} que é recebido numa \emph{potência} distingue-se
\emph{realmente} da \emph{potência} em que é recebido; porque repugna que uma
coisa seja recebida em si mesma. Ora, no \emph{ente finito}, a \emph{existência}
é recebida na \emph{essência} como o \emph{ato} na \emph{potência}; porquanto,
se a \emph{existência} não fosse recebida na \emph{essência} como o \emph{ato}
na \emph{potência}, a \emph{existência} dos \emph{entes finitos} seria um
\emph{ato puro}, \emph{subsistente}, e por isso seria \emph{infinito} e
\emph{único}: o que repugna. Logo, no \emph{ente finito}, a \emph{existência}
distingue-se \emph{realmente} da \emph{essência}.\footnote{Depois do que
deixamos dito neste e no precedente artigo, não é necessário insistir muito na
verdade \textemdash{} que o \emph{ser} (ou a \emph{existência}) dos \emph{entes
finitos}, se não fosse recebido na \emph{essência} como \emph{ato} na
\emph{potência}, seria um \emph{ato puro}, \emph{subsistente}, e por isso
\emph{infinito} e \emph{único}. Porquanto todo o \emph{ato} é \emph{limitado} e
\emph{multiplicado} pela \emph{potência} que o recebe; e por isso um \emph{ato}
que não é recebido na \emph{potência} não pode ser \emph{limitado} nem
\emph{multiplicado}, mas é necessariamente \emph{infinito} e \emph{único}. Ora
o \emph{ser}, ou a \emph{existência}, dos \emph{entes finitos} é um \emph{ato
essencialmente finito}, \emph{multiplicável} e \emph{multiplicado}.
\textemdash{} Nem se diga que o \emph{ato} não recebido na \emph{essência} pode
ser \emph{limitado} ou \emph{multiplicado} por uma \emph{causa extrínseca};
porque nenhuma coisa pode ser \emph{limitada} ou \emph{multiplicada} por uma
\emph{causa extrínseca} se é intrinsecamente \emph{ilimitada} e \emph{única}, e
o \emph{ato subsistente} não tem em si mesmo um \emph{princípio} que o possa
limitar ou multiplicar. Diz o \emph{Angélico} S. Tomás: «\textit{Esse
subsistens oportet esse infinitum, quia non terminatur aliquo recipiente\dots{}
Impossibile est quod sit duplex esse omnino infinitum. Esse enim quod est
infinitum omnem perfectionem essendi comprehendit; et sic, si duobus talis
adesset infinitas, non inveniretur quo unum ab altero differret}» (\emph{C.
Gent.}, l.~II, 52). \textemdash{} Este argumento pode resumir-se assim: no
\emph{ente finito}, a \emph{existência} refere-se à \emph{essência} como o
\emph{ato} à \emph{potência}. Ora o \emph{ato} distingue-se \emph{realmente} da
\emph{potência}. Logo, no \emph{ente finito}, a \emph{existência} distingue-se
\emph{realmente} da \emph{essência}.

  Em poucas palavras: \emph{não podem realmente identificar-se duas coisas que
  são opostas entre si}; aliás o \emph{ente} e o \emph{não-ente} seriam uma e a
  mesma coisa na realidade. Ora a \emph{essência} e o \emph{ser}, nas
  \emph{criaturas}, \emph{opõem-se por muitos modos}, como se opõem o
  \emph{necessário} e o \emph{contingente}, o \emph{eterno} e o
  \emph{temporâneo}, o \emph{limitante} e o \emph{limitado}, o \emph{divisor} e
  o \emph{diviso}, o \emph{não-causado} e o \emph{causado}. Logo, nas
  \emph{criaturas}, a \emph{essência} e o \emph{ser} distinguem-se
  \emph{realmente}.

  Nas obras do S. Doutor encontram-se muitos outros argumentos que provam a
  \emph{distinção real} entre a \emph{existência} e a \emph{essência} nos
  \emph{entes finitos} (cf.\ \emph{C. Gent.}, II, 52; \emph{Sum. Th.}, p.~1,
  q.~3, a.~4; q.~50, a.~2 ad~3; q.~75, a.~5 ad~4; q.~88, a.~2 ad~4; \emph{De
  Pot.}, q.~7, a.~2 ad~9; \emph{De Ver.}, q.~27, a.~1 ad~8; \emph{Com. in lib.
  De Hebd.}, l.~2). \textemdash{} Não insistimos, porque os argumentos aduzidos
  são suficientes para a \emph{demonstração da tese}.} \pg{318}

\medskip
\noindent Para a exata compreensão da nossa \emph{tese} e para a \emph{refutação
das dificuldades}, fazemos as seguintes considerações:

\medskip
\noindent\textbf{I.}\quad É certo que uma \emph{essência atual}, por isso mesmo
que é \emph{atual}, \emph{existe}. Mas a questão consiste em saber se a
\emph{atualidade} compete a essa \emph{essência} \emph{por si mesma}, enquanto
é \emph{essência}, ou em virtude de uma outra \emph{realidade}, \emph{realmente
distinta} dela, que é a \emph{existência}. Porquanto uma coisa é e diz-se
\emph{atual} ou porque é \emph{ato}, ou porque \emph{recebe o ato}. Ora a
\emph{essência atual} não é \emph{atual por si mesma}, como se ela fosse
\emph{ato}, ou se identificasse com o \emph{ato}, \textemdash{} mas é
\emph{atual} enquanto \emph{recebe o ato} que a completa na linha da realidade
e que é a \emph{existência}. Sendo \emph{atual} pelo fato de \emph{receber o
ato}, a \emph{essência} conserva sempre a \emph{razão} ou a \emph{natureza de
potência}, e por isso distingue-se \emph{realmente} da \emph{existência}, porque
o \emph{ato} e a \emph{potência} são \emph{coisas realmente distintas}.
\textemdash{} Se, às vezes, dizemos que a \emph{essência atual} é \emph{atual
por si mesma}, comparamos a \emph{essência atual} com a \emph{essência
possível}, e não com a \emph{existência}; porque, efetivamente, a \emph{essência
atual} refere-se à \emph{essência possível} como o \emph{ato} se refere à
\emph{potência}.

\medskip
\noindent\textbf{II.}\quad A \emph{essência atual}, ou \emph{física}, é dotada
de \emph{entidade própria}, e tem a sua \emph{última atualidade} na própria
linha \textemdash{} na \emph{linha da essência}, enquanto nada lhe falta do que
é denotado pela sua \emph{definição}. Mas esta \emph{última atualidade na linha
da essência} não é suficiente para explicar a sua \emph{presença na realidade}.
Diz S. Tomás: «\textit{Natura considerata praeter suum esse est secundum
acceptionem rationis tantum}» (\emph{In I Sent.}, dist.~14, q.~1, a.~3). Se,
pois, a \emph{essência} existe na realidade, existe em virtude de uma outra
\emph{atualidade} \textemdash{} da \emph{atualidade de uma ordem diversa}, que
é o \emph{ato da existência}. \pg{319} \textemdash{} A \emph{existência}, nos
\emph{entes finitos}, não é nem pode ser um \emph{elemento constitutivo} da
\emph{essência}; e por isso, quando se une com a \emph{essência}, não a
\emph{altera}, não a \emph{transforma}, deixa-a ficar o que era; limita-se a
\emph{completá-la na ordem da existência}, colocando-a e conservando-a colocada
na realidade.

\medskip
\noindent\textbf{III.}\quad A \emph{existência da criatura}, por isso mesmo que
é \emph{existência}, não é \emph{indiferente para o oposto}, isto é, para a
\emph{não-existência}; aliás uma coisa seria e deixaria de ser ao mesmo tempo e
sob o mesmo aspecto. Por isso, se uma coisa é \emph{indiferente} para a
\emph{existência} e para a \emph{não-existência}, esta \emph{indiferença} não
deriva da própria \emph{existência}, mas deriva do \emph{sujeito} em que a
\emph{existência} é recebida e que pode \emph{existir} e \emph{deixar de
existir}. \textemdash{} Como também, se a \emph{existência} acaba, não acaba
porque é \emph{existência}, mas porque acaba o seu \emph{sujeito}. A razão é
porque um \emph{ato misto}, como é a \emph{existência das criaturas}, segue a
sorte do \emph{sujeito} de que é \emph{ato}.

\medskip
\noindent\textbf{IV.}\quad O \emph{termo adequado da produção} não é só a
\emph{essência}, nem é só a \emph{existência}, mas é o \emph{composto de uma e
de outra}. Na verdade, a \emph{causa eficiente}, quando produz a \emph{essência}
e a \emph{existência}, não as produz como \emph{duas coisas}, mas como \emph{uma
coisa só}, composta de \emph{dois princípios}, um dos quais se refere ao outro
como a \emph{potência} ao \emph{ato}. \textemdash{} Por isso, nem a
\emph{existência} precede, na ordem do tempo, a \emph{essência}, nem a
\emph{essência} precede a \emph{existência}; mas, pelo fato mesmo de dar a
\emph{existência} à \emph{essência}, o \emph{agente} produz não só a
\emph{existência}, mas também a \emph{essência}. Diz S. Tomás: «\textit{Ex hoc
ipso quod quidditati esse tribuitur, non solum esse sed ipsa quidditas creari
dicitur; quia antequam esse habeat, nihil est, nisi forte in intellectu
creantis, ubi non est creatura, sed creatrix essentia}» (\emph{De Pot.}, q.~3,
a.~5 ad~2).

\medskip
\noindent\textbf{V.}\quad Donde se segue que a \emph{existência} e a
\emph{essência}, embora sejam \emph{coisas realmente distintas}, contudo não se
podem separar uma da outra sem que uma e outra imediatamente acabem. Porquanto
o \emph{efeito} que de um \emph{ato} deriva para um \emph{sujeito} não permanece
e não pode permanecer se acabar o próprio \emph{ato}. Ora \emph{existir} é o
\emph{efeito próprio da existência}, e esse \emph{efeito} não pode ser produzido
nem executado senão na \emph{essência}. Repugna, pois, que a \emph{essência}
dure sem a \emph{existência}; como repugna que a \emph{existência} subsista sem
a \emph{essência}, para a qual tem uma \emph{relação necessária}. \textemdash{}
Nunca devemos esquecer que o \emph{ente existente} é o \emph{composto}, e que a
\emph{existência} e a \emph{essência} não são \emph{duas realidades
subsistentes}, mas são os \emph{princípios intrínsecos} \textemdash{}
\emph{atual} e \emph{potencial} \textemdash{} pelos quais o \emph{composto}
existe.

\teseB{45}{No Ente infinito, a existência não se distingue realmente da
essência}

\subp{a} Uma coisa não se distingue \emph{realmente} de outra quando entre elas
há \emph{perfeita identidade}. Ora, no \emph{Ente infinito}, há \emph{perfeita
identidade} entre a \emph{existência} e a \emph{essência}; porque Deus existe em
virtude ou pela \emph{exigência} da sua \emph{essência}, e por isso a
\emph{existência} constitui a sua \emph{essência}. Logo, no \emph{Ente
infinito}, a \emph{existência} não se distingue \emph{realmente} da
\emph{essência}.

\subp{b} Se, no \emph{Ente infinito}, a \emph{existência} se distinguisse
\emph{realmente} da \emph{essência}, a \emph{existência} seria recebida na
\emph{essência}, seria \emph{ato misto}, \emph{não-subsistente}, e por isso
\emph{limitado}, e \emph{multiplicado} pela mesma \emph{essência} em que é
recebido. Ora a \emph{existência}, ou o \emph{ser de Deus}, é \emph{ato
puríssimo}, \emph{subsistente}, e por isso \emph{infinito} e \emph{único na sua
ordem}. Logo, no \emph{Ente infinito}, a \emph{existência} não se distingue
\emph{realmente} da \emph{essência}.\footnote{Deus, e só Ele, é o \emph{ente
necessário} que existe em virtude da sua \emph{essência}, isto é, por assim o
exigir a sua \emph{essência}; e por isso em Deus, e só nEle, a \emph{existência}
não se distingue \emph{realmente} da \emph{essência}, mas identifica-se com
esta. De fato, todos dizem \textemdash{} que Deus, e só Ele, é, \textemdash{}
que Deus é o \emph{próprio Ser}, a \emph{própria Existência}, \textemdash{} que
todos os \emph{seres finitos} não são o \emph{ser}, mas receberam o \emph{ser},
ou a \emph{existência}.

  Donde se vê que a \emph{essência} pode \emph{atuar-se} ou \emph{realizar-se}
  por dois modos: \emph{por si mesma}, ou \emph{pela ação de uma causa
  extrínseca}. \emph{Realiza-se por si mesma} quando é \emph{necessária}, quando
  existe na realidade por \emph{virtude própria}, por assim ela o exigir; assim
  existe a \emph{Essência de Deus}. \emph{Realiza-se pela ação de uma causa
  extrínseca} quando é \emph{contingente}, isto é, \emph{indiferente para a
  existência e para a não-existência}, e por isso deve ser determinada por uma
  \emph{causa extrínseca}; assim existe a \emph{essência do ente finito}.
  \textemdash{} Quando a \emph{essência} se \emph{realiza por si mesma}, a
  \emph{existência} é um \emph{ato constitutivo} da mesma \emph{essência}, isto
  é, a \emph{essência} não só \emph{existe}, mas \emph{é a sua existência}, o
  \emph{seu ser}. \pg{320} Quando a \emph{essência} é determinada para a
  \emph{existência} pela \emph{ação de uma causa extrínseca}, a \emph{existência}
  não é um \emph{elemento constitutivo} da \emph{essência}, mas é uma
  \emph{perfeição} que se lhe une, como o \emph{ato} se une com a
  \emph{potência}, e que lhe dá o \emph{último complemento}, a \emph{última
  atualidade} na linha da realidade.

  Tal é a \emph{doutrina de S. Tomás}. O S. Doutor não se cansa de repetir
  \textemdash{} que em Deus o \emph{ser} e a \emph{essência} são \emph{uma e a
  mesma coisa}, \textemdash{} que Deus é \emph{ato subsistente}, \emph{forma
  puríssima}, \emph{não recebida em nenhuma potência}, \textemdash{} e que por
  isso Deus é \emph{infinito na sua perfeição} e \emph{único na sua ordem}, ao
  passo que, nas \emph{criaturas}, o \emph{ser} é uma \emph{coisa diversa da
  essência}, o \emph{ser} é \emph{recebido na essência}, e por isso as
  \emph{criaturas} são \emph{finitas} e \emph{muitas}. Eis as suas palavras:
  «\textit{Esse est actualitas omnis formae vel naturae; non enim bonitas vel
  humanitas significatur in actu, nisi prout significamus eam esse. Oportet
  igitur quod ipsum esse comparetur ad essentiam, quae est aliud ab ipso, sicut
  actus ad potentiam. Cum igitur in Deo nihil sit potentiale, sequitur quod non
  sit aliud in eo essentia quam suum esse. Sua igitur essentia est suum esse}»
  (\emph{Sum. Theol.}, p.~1, q.~3, a.~4). \textemdash{} Mais adiante:
  «\textit{Infinitum quod se tenet ex parte formae, non determinata per
  materiam, habet rationem perfecti\dots{} Illud autem quod est maxime formale
  omnium est ipsum esse. Cum igitur esse divinum non sit esse receptum in
  aliquo, sed ipse sit suum esse subsistens, manifestum est quod ipse Deus sit
  infinitus et perfectus}» (\emph{Sum. Th.}, p.~1, q.~7, a.~1). \textemdash{}
  Noutro lugar diz: «\textit{Talis res, quae sit suum esse, non potest esse
  nisi una; unde oportet quod in qualibet alia re, praeter eam, sit aliud suum
  esse et aliud quidditas}» (\emph{De Ente et Essentia}, c.~V). \textemdash{}
  «\textit{Huiusmodi formae (Angeli) non terminantur neque contrahuntur per
  aliquam materiam. Sed quia forma creata sic subsistens habet esse, et non est
  suum esse, necesse est quod ipsum eius esse sit receptum et contractum ad
  terminatam naturam. Unde non potest esse infinitum simpliciter}» (\emph{Sum.
  Th.}, p.~1, q.~7, a.~2). \textemdash{} «\textit{Esse hominis terminatum est ad
  hominis speciem, quia est receptum in natura speciei humanae. Esse autem Dei,
  cum non sit in aliquo receptum, sed sit esse purum, non limitatur ad aliquem
  modum perfectionis}» (\emph{De Pot.}, q.~1, a.~2; cf.\ \emph{De spirit.
  creat.}, a.~1). \textemdash{} Nos citados trechos, que se referem aos
  \emph{entes finitos} ou \emph{criados}, trata-se de uma \emph{distinção real}
  entre a \emph{essência} e a \emph{existência}; porque, se se tratasse apenas
  de uma \emph{distinção lógica}, não seria legítima a argumentação do S.
  Doutor, visto que também em Deus há uma \emph{distinção lógica} entre a
  \emph{essência} e a \emph{existência}. \textemdash{} Mas, deixando outras
  citações (cf.\ \emph{I Sent.}, dist.~19, q.~2, a.~2; \emph{C. Gentes}, l.~II,
  c.~52), o trecho seguinte torna a coisa evidente. No \emph{Comm.} ao livro de
  Boécio \emph{De Hebdomad.} escreve: «\textit{Est ergo considerandum quod
  sicut esse et quod est differunt in simplicibus secundum intentionem, ita in
  compositis differunt realiter}» (Lect.~2). O que se entende por \emph{ente
  simples} é dito com as seguintes palavras: «\textit{Hoc autem simplex unum et
  sublime est ipse Deus}». Poderia ser mais claro? Em Deus, que é o \emph{único
  ente simplicíssimo}, a \emph{essência} e a \emph{existência} diferem
  \emph{logicamente}; em todas as \emph{criaturas}, que são \emph{entes
  compostos}, \pg{321} a \emph{essência} e a \emph{existência} diferem
  \emph{realmente} (cf.\ \emph{De Verit.}, q.~27, a.~1 ad~8).

  A \emph{distinção real}, pois, entre a \emph{essência} e a \emph{existência},
  nas \emph{criaturas}, é doutrina certamente propugnada pelo \emph{Angélico
  Mestre} e é um dos \emph{pontos capitais e fundamentais} de toda a sua
  \emph{metafísica}, como afirmam, entre outros, o \emph{Card.\ Pallavicino}
  (\emph{De Deo}, c.~III) e \emph{Silvestre Mauro} (\emph{Quaest. phil.},
  vol.~II, q.~7). \textemdash{} Poderá impugnar-se esta doutrina; mas não pode,
  sem grande \emph{temeridade}, afirmar-se que é falso ou incerto que ela tenha
  sido ensinada por S. Tomás. \emph{Tifano}, que todavia não segue o
  \emph{Angélico Mestre} neste ponto, escreve: «\textit{Quod existentiam in
  creaturis re distinguat (S. Thomas) ab essentia, id vel negare, vel in dubium
  revocare, est hominis impudentis, vel in eius doctrina peregrini}» (\emph{De
  hypostasi}, c.~6).}

\tese{46}{Importância da doutrina exposta} A \emph{doutrina exposta} tem uma
grande \emph{importância}, porque serve para estabelecer uma \emph{harmonia
perfeitíssima} entre todos os seres. A \emph{simplicidade} é \emph{sinal de
perfeição}, como a \emph{composição} é \emph{sinal de limitação ou imperfeição}.
Deus é \emph{ato simplicíssimo}, \emph{puríssimo}; a sua \emph{existência}
identifica-se com a \emph{essência}; por isso é \emph{infinito na perfeição}.
\textemdash{} As \emph{inteligências separadas}, que são os \emph{Anjos}, são
compostas de \emph{essência} e de \emph{existência}; por isso são \emph{limitadas
na perfeição}, infinitamente inferiores a Deus, mas superiores a todas as outras
\emph{criaturas}. \textemdash{} As \emph{criaturas visíveis} são compostas de
\emph{essência} e de \emph{existência}, e também de \emph{matéria} e de
\emph{forma}; por isso são ainda mais \emph{limitadas}, e \emph{ínfimo} é o seu
grau de perfeição.\footnote{Nem essa é a única vantagem. A questão acerca da
\emph{distinção real} entre a \emph{existência} e a \emph{essência}, nas
\emph{criaturas}, liga-se intimamente com as outras duas importantes questões
acerca da \emph{distinção real} entre a \emph{substância} e os \emph{acidentes},
e entre a \emph{matéria} e a \emph{forma}. Estas três questões refletem, pela
sua vez, muita luz sobre alguns \emph{dogmas da Fé}. Na verdade, a \emph{teoria
da distinção real} entre a \emph{essência} e a \emph{existência} explica, de
algum modo, como é que, em N.~S. Jesus Cristo, podem \emph{duas naturezas} ou
\emph{essências} subsistir numa só \emph{existência}, ou \emph{pessoa divina};
\textemdash{} a \emph{teoria da distinção real} entre a \emph{substância} e os
\emph{acidentes} serve para conhecer, de algum modo, o \emph{mistério da
Eucaristia}; \textemdash{} a \emph{teoria da matéria e da forma} explica a
\emph{união da alma com o corpo} em harmonia com as \emph{definições da Igreja}.}

\end{document}
